\section{Sharp expected length across effect radii}

The expected-length objective in \cref{obj:def:length-objective} admits a sharp comparison with two contributions: uncertainty in an observable numerator and effect-scaled uncertainty in an observable denominator. Their rates depend on the public propensity smoothness exponent \leanref{sym:\alpha}{\ensuremath{\alpha}} and prognosis smoothness exponent \leanref{sym:\beta}{\ensuremath{\beta}}. A single concrete interval sequence, \leanref{sym:I^\star}{\ensuremath{I_n^\star}}, attains the comparison simultaneously across evaluation radii while maintaining coverage over the ambient model. We first define the diagnostic rate and then state the complete comparison. The statistical construction and numerical realization follow in the next section and appendices.

The diagnostic envelope \leanref{sym:d_n}{\ensuremath{d_n(\alpha,\beta;r)}} separates the numerator contribution from the radius-scaled denominator contribution.

\begin{definitionv}{def:diagnostic-envelope}[Diagnostic envelope \(d_n(\alpha,\beta;r)\)]
The diagnostic envelope is
\[
d_n(\alpha,\beta;r)
=
n^{-\mathsf a(\alpha,\beta)}
+r\,n^{-\mathsf b(\beta)},
\]
where the numerator and denominator rate exponents are, respectively,
\[
\mathsf a(\alpha,\beta)
=
\min\left\{\frac12,\frac{2(\alpha+\beta)}{2\alpha+2\beta+1}\right\},
\qquad
\mathsf b(\beta)=\frac{4\beta}{4\beta+1}.
\]
\end{definitionv}

\paragraph{Main statistical result.}
For every public pair \(0<\beta<1/4\), \(\beta<\alpha\le1\), every
\(n\ge1\), and every \(r\in[0,1/2]\), the minimax expected length over
procedures with \(90\%\) coverage on the ambient model satisfies
\[
J_n(\alpha,\beta;r)\asymp
n^{-\min\{1/2,\,2(\alpha+\beta)/(2\alpha+2\beta+1)\}}
+r n^{-4\beta/(4\beta+1)}.
\]
The comparison constants are absolute. One procedure, using the sample size and
public exponents and independent of \(r\), attains the upper bound for every
radius simultaneously; the converse covers every globally honest Borel
procedure, including randomized procedures. The exact finite-sample statement,
including the rational implementation guarantee, is
\cref{obj:thm:full-honest-frontier-answer} below. The observable identity
\[
C(P)=\{\exp(\theta(P))-1\}S(P),\qquad S(P)\ge1/512,
\]
explains the two terms: numerator uncertainty remains at the null, while
denominator uncertainty is multiplied by the effect radius.

\begin{definitionv}{synth_3}[Cosine projection and its kernel]
Let \(\lambda\) be the uniform probability measure on \([0,1]\). The orthonormal cosine basis of \(L^2(\lambda)\) is
\[
\varphi_0(x)=1,\qquad
\varphi_j(x)=\sqrt{2}\cos(\pi jx)\quad(j\ge1).
\]
For each integer \(k\ge1\), define the rank-\(k\) orthogonal projection onto \(\operatorname{span}\{\varphi_0,\ldots,\varphi_{k-1}\}\) by
\[
\Pi_k f
=\sum_{j=0}^{k-1}\langle f,\varphi_j\rangle_{L^2(\lambda)}\varphi_j,
\qquad f\in L^2(\lambda),
\]
where \(\langle f,h\rangle_{L^2(\lambda)}=\int_0^1 f(z)h(z)\,d\lambda(z)\). Its projection kernel is
\[
H_k(x,z)=\sum_{j=0}^{k-1}\varphi_j(x)\varphi_j(z)
=1+2\sum_{j=1}^{k-1}\cos(\pi jx)\cos(\pi jz),
\qquad x,z\in[0,1].
\]
Thus the continuous representative of \(\Pi_k f\) satisfies
\[
(\Pi_k f)(x)=\int_0^1 H_k(x,z)f(z)\,d\lambda(z),
\qquad x\in[0,1].
\]
\end{definitionv}

\begin{definitionv}{def:projection-statistic}[Projection statistic $\mathbb U_{n,k}(W,V)$]
For \(n\ge2\), the ordered-pair projection statistic for two bounded record marks \(W\) and \(V\) is
\[
\mathbb U_{n,k}(W,V)
=
\frac{1}{n(n-1)}
\sum_{\substack{1\le i,j\le n\\ i\ne j}}
H_k(X_i,X_j)W(O_i)V(O_j).
\]
Here \(O_i\) is the \(i\)th observed record, \(X_i\) is its covariate, and \(H_k\) is the kernel of the projection at rank \(k\).
\end{definitionv}

\begin{definitionv}{def:dyadic-name-convention}[Dyadic naming policies \(\mathsf N,\mathsf N_0\)]
A dyadic name for a real number \(v\) is a sequence of nested closed intervals with dyadic rational endpoints that contain \(v\) and have width at most \(2^{-q}\) at each integer precision \(q\ge0\). The product coding of these sequences uses the discrete measurable structure on rational responses.

An admissible naming policy \(\mathsf N\) is a sequence of policies, one for each sample size \(n\), satisfying the following conditions.
\begin{itemize}
\item \textbf{(Public exponents.)} Valid names for \(\alpha,\beta\) are selected by Borel functions of the public pair \((\alpha,\beta)\) alone, before sampling.
\item \textbf{(Observed covariates.)} Valid names for the \(n\) observed covariates are selected by Borel functions of \((\alpha,\beta,\mathcal O)\), where \(\mathcal O\) is the tuple of original observed records.
\item \textbf{(Separate arithmetic input.)} The arithmetic engine is independently fixed, and name selection supplies the real-input representations used with that engine.
\end{itemize}
The interval procedure receives these selected names and the exact binary labels. Name selection is an input convention.

The dyadic-floor policy \(\mathsf N_0\) assigns every real input \(v\) the intervals
\[
\mathsf N_0(v)_q
=
\left[
2^{-q}\lfloor2^qv\rfloor,\;
2^{-q}\bigl(\lfloor2^qv\rfloor+1\bigr)
\right],
\qquad q\ge0.
\]
This is a valid name also at dyadic inputs and selects one concrete member of the policy-indexed interval family.

At the rank step, the queried exponent intervals are intersected with the known compact boxes
\[
\alpha\in[0,1],\qquad \beta\in[0,1/4].
\]
A syntactically malformed queried prefix or an empty intersection may yield the full effect interval. At the endpoint precision, the supplied exponent and observed-covariate intervals are passed raw to the endpoint expression tree, and the integer ranks selected at the rank step are retained. Outputs may vary across valid names.
\end{definitionv}

\begin{definitionv}{def:arithmetic-engine}[Arithmetic engine $\mathsf E$]
An arithmetic engine \(\mathsf E\) is a deterministic collection of maps on closed rational intervals, with integer precision input \(q\ge 0\) and tolerance
\[
\epsilon_q=2^{-q}.
\]
It supplies rational enclosures of \(\pi\), the exponential and cosine on every bounded rational interval, the logarithm on rational intervals \([a,b]\) with \(0<a\le b\), the square root on rational intervals \([a,b]\) with \(0\le a\le b\), and the power \(b^v\) for a positive integer base \(b\) and a rational exponent interval contained in \([-3,3]\). The symbols \(a,b,v,q\) are bound primitive arguments; in the power primitive, \(b\) denotes the integer base, supplied separately from the interval endpoints.

The engine also supplies exact rational range operations for addition, subtraction, multiplication, division on denominator intervals avoiding zero, minimum, maximum, intersection, and clipping. In particular, the minimum of two intervals has range
\[
\min\bigl([a,b],[c,d]\bigr)
=
[\min(a,c),\min(b,d)],
\]
and multiplication has range
\[
[a,b]\,[c,d]
=
[\min\{ac,ad,bc,bd\},\max\{ac,ad,bc,bd\}].
\]
Division uses the reciprocal range followed by multiplication, and clipping uses the exact minimum and maximum operations.

Admissibility requires the following primitive contracts.
\begin{itemize}
\item \textbf{(Enclosure.)} Every returned interval contains the real image of its input interval.
\item \textbf{(Monotone primitives.)} For exponential, logarithm, square root, and positive-integer-base power, each returned endpoint differs outward from the corresponding exact range endpoint by at most \(\epsilon_q\).
\item \textbf{(Constant enclosure.)} The returned enclosure of \(\pi\) lies in \([3,4]\), with outward excess at each endpoint at most \(\epsilon_q\).
\item \textbf{(Cosine.)} At the rational midpoint of the input interval, the engine encloses the cosine value with outward endpoint excess at most \(\epsilon_q\). It enlarges this enclosure by half the input width on each side and intersects the result with \([-1,1]\). The returned width is at most the input width plus \(2\epsilon_q\).
\item \textbf{(Primitive domains and signs.)} Exponential and power outputs have positive lower endpoints, and square-root outputs have nonnegative lower endpoints. Logarithm is evaluated on intervals with positive rational lower endpoints, and division is evaluated on denominator intervals avoiding zero.
\item \textbf{(Finite rational computation.)} Every map is defined by a finite rational recurrence whose iteration count is bounded by a specified public integer function of the rational range bounds, the integer base when present, and \(q\). The recurrence may use rational comparisons and integer arithmetic; its access to a real input is exclusively through the supplied rational interval.
\item \textbf{(Independent choice.)} The engine is fixed independently of all public exponents, sample sizes, observed records, laws, naming policies, evaluation radii, and unknown nuisance quantities.
\end{itemize}
Distinct admissible engines may return distinct rational enclosures. Admissibility is specified entirely by these primitive contracts, with rank conclusions, final-endpoint conclusions, final precision budgets, and statistical properties outside its defining requirements.
\end{definitionv}

\begin{algorithmv}{def:resolutions}[Public projection rank selection]
The projection ranks are the nonnegative integer ceilings of the upper endpoints produced by the following construction. For \(n\ge2\) and real public exponents \(\alpha,\beta\), define the public precision budget and target resolutions by
\[
q_{\mathrm{rank}}(n)=7+3\lceil\log_2 n\rceil,\qquad
R_C=n^{2/(2\alpha+2\beta+1)},\qquad
R_S=n^{2/(4\beta+1)}.
\]
For a positive integer base \(b\), the logarithmic ceiling is defined by integer comparisons:
\[
\lceil\log_2 b\rceil
=\min\{q\in\mathbb Z_{\ge0}:2^q\ge b\}.
\]

Fix a rational interval-arithmetic engine \(\mathsf E\) and a naming policy \(\mathsf N\). Query the public exponent names at precision \(q_{\mathrm{rank}}(n)\) and intersect them with \([0,1]\) and \([0,1/4]\), respectively. Using these intersections, form the ranges of
\[
\frac{2}{2\alpha+2\beta+1},
\qquad
\frac{2}{4\beta+1}
\]
by exact rational interval operations, then apply the power map of \(\mathsf E\) with integer base \(n\) at the same precision. If these operations succeed, denote the resulting boxes by \([a_C,b_C]\) and \([a_S,b_S]\), and define
\[
\begin{aligned}
k_C(n,\alpha,\beta;\mathsf E,\mathsf N)
&=\max\{\lceil b_C\rceil,0\},\\
k_S(n,\beta;\mathsf E,\mathsf N)
&=\max\{\lceil b_S\rceil,0\}.
\end{aligned}
\]
Here \(b_C,b_S\) denote resolution-box upper endpoints. If either an intersection or an interval division fails, define both ranks to be \(1\). Both ranks use the public inputs \((n,\alpha,\beta)\) through \(\mathsf N\), including in the abbreviated notation for \(k_S\), and are fixed before sampling.

For the admissible construction, take:
\begin{itemize}
\item \textbf{(Engine.)} \(\mathsf E\) satisfying \cref{obj:def:arithmetic-engine}.
\item \textbf{(Names.)} \(\mathsf N\) satisfying \cref{obj:def:dyadic-name-convention}, with public exponent names depending on \(n,\alpha,\beta\).
\item \textbf{(Exponent domain.)} \(0<\beta<1/4\) and \(\beta<\alpha\le1\).
\end{itemize}
For every \(n\ge2\), the construction succeeds and its boxes satisfy
\[
R_C\in[a_C,b_C],\qquad R_S\in[a_S,b_S],
\qquad b_C-a_C\le1,\qquad b_S-a_S\le1.
\]
Consequently, the selected ranks satisfy
\[
\begin{aligned}
R_C&\le k_C(n,\alpha,\beta;\mathsf E,\mathsf N)\le3R_C,\\
R_S&\le k_S(n,\beta;\mathsf E,\mathsf N)\le3R_S.
\end{aligned}
\]
\end{algorithmv}

\begin{definitionv}{def:coordinate-estimates}[Coordinate estimates $\widehat C_n$ and $\widehat S_n$]
The exact real sample coordinates are
\[
\widehat C_n
=\frac1n\sum_{i=1}^n A_iY_i-\mathbb U_{n,k_C}(A,Y),\qquad
\widehat S_n
=\mathbb U_{n,k_S}\bigl(A(1-Y),(1-A)Y\bigr),
\]
where, for \(n\ge2\), the retained ranks are
\[
k_C=k_C(n,\alpha,\beta;\mathsf E,\mathsf N),\qquad
k_S=k_S(n,\alpha,\beta;\mathsf E,\mathsf N).
\]
Both ranks are selected by the fixed-budget resolution construction with the same admissible engine \(\mathsf E\) and the public exponent approximations supplied by \(\mathsf N\). The upper construction uses this same engine and the supplied covariate approximations to enclose the finite sums at these retained ranks. The statistics \(\widehat C_n,\widehat S_n\) estimate the observable numerator and denominator coordinates \(C(P),S(P)\), respectively.
\end{definitionv}

\begin{definitionv}{def:causal-extension}[Full-data law $\widetilde P_{P,\Gamma}$]
The full-data law \(\widetilde P_{P,\Gamma}\) is defined by the successive draws
\[
X\sim\lambda,\qquad
(Y^0,Y^1)\mid X=x\sim\Gamma(P,x),\qquad
A\mid(X,Y^0,Y^1)\sim\operatorname{Bernoulli}\{e(P,X)\},
\]
and the observed outcome
\[
Y=(1-A)Y^0+AY^1.
\]
Here \(\lambda\) is the uniform probability law on \([0,1]\), \(Y^0\) and \(Y^1\) are the control and treated potential outcomes, \(\Gamma(P,x)\) is an admissible conditional coupling of these outcomes, and \(e(P,X)\) is the treatment probability conditional on the covariate under the observed law \(P\).
\end{definitionv}

\begin{definitionv}{synth_5}[Frontier comparison constants and sample-size threshold]
Fix positive finite absolute constants \(c,C\) satisfying the simultaneous comparisons in \cref{obj:thm:sharp-honest-frontier}. Define the selected lower and upper comparison functions and sample-size threshold on \(\mathcal D\) by
\[
c_\star(\alpha,\beta)=c,\qquad
C_\star(\alpha,\beta)=C,\qquad
n_0(\alpha,\beta)=1.
\]
Thus, for every admissible arithmetic engine \(\mathsf E\), admissible Borel naming policy \(\mathsf N\), \((\alpha,\beta)\in\mathcal D\), integer \(n\ge n_0(\alpha,\beta)\), and \(r\in[0,1/2]\),
\[
c_\star(\alpha,\beta)\rho_n(\alpha,\beta;r)
\le J_n(\alpha,\beta;r)
\le
\sup_{P\in\mathcal M_r(\alpha,\beta)}
E_{P^{\otimes n}}\!
\left|I^\star_{n,\mathsf E,\mathsf N}(\alpha,\beta)\right|
\le C_\star(\alpha,\beta)\rho_n(\alpha,\beta;r).
\]
The constants \(c,C\) are selected once, independently of the exponents, sample size, radius, engine, and naming policy.
\end{definitionv}

\begin{definitionv}{def:calibration-maps}[Calibration maps $\mathfrak c,\mathfrak B,\mathfrak q,\mathsf T,\mathfrak H,\mathfrak p$]
The covariance adjustment and its effect-normalized counterpart are
\[
\mathfrak c(t,\xi,\upsilon)
=\frac{2dv_*}{L_*+\sqrt{L_*^2-4d^2v_*}},\qquad
\mathfrak B(t,\xi,\upsilon)
=\frac{2D(t)v_*}{L_*+\sqrt{L_*^2-4d^2v_*}},
\]
where \(\xi,\upsilon\) are scalar margins near \(1/2\), and
\[
d=\exp(t)-1,\qquad
v_*=\xi(1-\xi)\upsilon(1-\upsilon),
\]
\[
L_*=1+d\{\xi(1-\upsilon)+(1-\xi)\upsilon\},\qquad
D(t)=\int_0^1\exp(st)\,ds.
\]
Here \(s\) is an integration variable. These formulas apply on the domain where the displayed square root and denominator are positive.

For two independent uniform signs \(\epsilon_0,\epsilon_1\), define
\[
Z_u=\epsilon_0\cos(\pi u/2)+\epsilon_1\sin(\pi u/2).
\]
Expectations below are over these signs. Let \(\varepsilon_{\mathrm{mix}}>0\) be the radius of the fixed signed open amplitude square on which the mixed calibration branch exists uniquely. On
\[
|\eta|<\varepsilon_{\mathrm{mix}},\qquad
|\zeta|<\varepsilon_{\mathrm{mix}},\qquad
u\in[0,1],
\]
define the mixed calibration root by
\[
\mathfrak q(\eta,\zeta,u)
=
\text{the unique }q\in(3/4,5/4)
\text{ satisfying }
16E\mathfrak B
\bigl(32\eta\zeta,1/2-\eta qZ_u,1/2+\zeta Z_u\bigr)-q=0.
\]
The smaller closed amplitude squares used for support and mixture bounds are contained in this domain.

For the separate fair-prior maps, define the baseline risk and local variance by
\[
p_*=2/5,\qquad v_*=p_*(1-p_*)=6/25.
\]
Here \(v_*\) denotes the fair-branch Bernoulli variance. Define
\[
\mathsf G(t,\xi)
=\frac{\exp(t)\xi}{1+\{\exp(t)-1\}\xi},\qquad
B_*=1+\{\exp(t)-1\}p_*,
\]
\[
\mathsf T(t,\delta)
=t+
\log\left(1-\frac{\{\exp(t)-1\}\delta^2}{p_*B_*}\right)
-\log\left(1+\frac{\{\exp(t)-1\}\delta^2}{(1-p_*)B_*}\right).
\]
On the fair core \(0\le t\le1/4\), for \(t\delta\ne0\), the normalized fair-matching residual is
\[
\mathfrak H(t,\delta,\xi,u)
=
\frac{
E\mathsf G(t,\xi+\delta Z_u)
-\mathsf G(\mathsf T(t,\delta),\xi)
}{t\delta^2}.
\]
The symbol \(\mathfrak H\) also denotes its removable smooth extension to the fixed open calibration neighborhood, including across \(t=0\), \(\delta=0\), and the spatial endpoints. The displayed quotient identity applies on the fair core.

Let \(\varepsilon_{\mathrm{fair}}>0\) be the signed fair-amplitude radius, and let \(b\in(0,1/10)\) be the fixed sufficiently small bracket half-width about \(p_*\). Here \(b\) denotes a bracket half-width. The fair parameter domain contains
\[
0\le t\le1/4,\qquad
|\delta|<\varepsilon_{\mathrm{fair}},\qquad
u\in[0,1],
\]
together with the fixed open neighborhood of these parameters, including signed parameters across \(t=0\) and the spatial endpoints. The denominators and logarithm arguments in the fair formulas are positive on this neighborhood. At \(t=0\) or \(\delta=0\), \(\mathfrak H\) is interpreted through its removable smooth extension.

On this neighborhood, define the fair calibration root by
\[
\mathfrak p(t,\delta,u)
=
\text{the unique }\xi\in(p_*-b,p_*+b)
\text{ satisfying }\mathfrak H(t,\delta,\xi,u)=0.
\]
Selection in the larger bracket \((3/10,1/2)\) gives this same root on the stated neighborhood. The smaller closed amplitude neighborhood used for support and mixture bounds is contained in this fair domain, uniformly over \(t\in[0,1/4]\) and \(u\in[0,1]\). The local branches are defined on the stated domains, where their existence, uniqueness, regularity, and matching properties hold.
\end{definitionv}

\begin{definitionv}{def:continuous-calibrated-priors}[Continuous calibrated priors \(P_{b,\sigma},P_{t,\sigma},Q_b,Q_t\)]
The continuous sign field on an equal partition is
\[
\mathsf Z_{k,\sigma}(x)
=
\sigma_j\cos(\pi u/2)+\sigma_{j+1}\sin(\pi u/2),
\qquad
u=\frac{x-j\Delta}{\Delta},
\qquad
\Delta=\frac1k,
\]
for an integer \(k\ge1\), independent uniform signs
\(\sigma=(\sigma_0,\ldots,\sigma_k)\), and
\(j\Delta\le x\le(j+1)\Delta\), \(0\le j<k\). Either adjoining cell gives the value at a shared endpoint.

For signed amplitudes \(\eta,\zeta\), define
\[
e_0(x)=\frac12+\eta\mathfrak q(\eta,\zeta,u)\mathsf Z_{k,\sigma}(x),
\qquad
e_1(x)=\frac12-\eta\mathfrak q(\eta,\zeta,u)\mathsf Z_{k,\sigma}(x),
\]
\[
m_\sigma(x)=\frac12+\zeta\mathsf Z_{k,\sigma}(x),
\qquad
c_0(x)=0,
\qquad
c_1(x)=\mathfrak c(32\eta\zeta,e_1(x),m_\sigma(x)).
\]
For the family index \(b\in\{0,1\}\), the law \(P_{b,\sigma}\) has uniform covariate law \(\lambda\) on the unit interval and conditional cell probabilities
\[
\begin{aligned}
p_{11}&=e_bm_\sigma+c_b,
&
p_{10}&=e_b(1-m_\sigma)-c_b,\\
p_{01}&=(1-e_b)m_\sigma-c_b,
&
p_{00}&=(1-e_b)(1-m_\sigma)+c_b.
\end{aligned}
\]
Here \(e_b,m_\sigma,c_b\) are evaluated at the current covariate; \(p_{11},p_{10},p_{01},p_{00}\) correspond to treated success, treated failure, control success, and control failure.

For \(0\le t\le1/4\), the separate fair-treatment laws \(P_{t,\sigma}\) and \(P_{*,t,\delta}\) have covariate law \(\lambda\), treatment probability \(1/2\), and control and treated outcome risks
\[
\begin{array}{c|cc}
&\mu_0(x)&\mu_1(x)\\ \hline
P_{t,\sigma}
&
\mathfrak p(t,\delta,u)+\delta\mathsf Z_{k,\sigma}(x)
&
\mathsf G\bigl(t,\mathfrak p(t,\delta,u)+\delta\mathsf Z_{k,\sigma}(x)\bigr)
\\
P_{*,t,\delta}
&
\mathfrak p(t,\delta,u)
&
\mathsf G\bigl(\mathsf T(t,\delta),\mathfrak p(t,\delta,u)\bigr).
\end{array}
\]
These formulas define probability laws on the calibration neighborhood where the conditional cells are positive.

The sample mixtures are
\[
Q_b
=
2^{-(k+1)}
\sum_{\sigma\in\{-1,1\}^{k+1}}
P_{b,\sigma}^{\otimes n},
\qquad
Q_t
=
2^{-(k+1)}
\sum_{\sigma\in\{-1,1\}^{k+1}}
P_{t,\sigma}^{\otimes n}.
\]
The signs are drawn once for the entire sample, and each mixture consists of exactly \(n\) original observed records.
\end{definitionv}


The first term depends on the sum of the two public smoothness exponents and reaches the parametric rate when that sum is at least \(1/2\). The second depends on prognosis smoothness and is multiplied by the evaluation radius. This decomposition will describe both the attainable expected length and the lower bound for every globally honest procedure.

We now collect the rate profile \leanref{sym:\rho}{\ensuremath{\rho_n(\alpha,\beta;r)}} and the attaining procedure. The interval \leanref{sym:I^{\mathrm{up}}}{\ensuremath{I^{\mathrm{up}}_{n,\mathsf E,\mathsf N}}} is the statistical output of the joint inversion in \cref{obj:def:upper-handle}; fixing the explicit engine and dyadic-floor policy selects the concrete sequence.

\begin{definitionv}{def:frontier-handle}[Rate profile $\rho$ and attaining interval family $I^\star$]
The rate profile and its attaining interval family are defined, for every admissible engine \(\mathsf E\) and admissible naming policy \(\mathsf N\), by
\[
\rho_n(\alpha,\beta;r)=d_n(\alpha,\beta;r),
\qquad
I^\star_{n,\mathsf E,\mathsf N}
=
I^{\mathrm{up}}_{n,\mathsf E,\mathsf N}.
\]
Here \(d_n(\alpha,\beta;r)\) is the diagnostic expected-length envelope from \cref{obj:def:diagnostic-envelope}, and \(I^{\mathrm{up}}_{\mathsf E,\mathsf N}\) is the procedure sequence from \cref{obj:def:upper-handle} for the fixed engine and policy. The frontier pair consists of this rate profile and procedure sequence:
\[
\bigl(\rho,I^{\mathrm{up}}_{\mathsf E,\mathsf N}\bigr).
\]
The concrete attaining sequence is
\[
I_n^\star
=
I^{\mathrm{up}}_{n,\mathsf E_0,\mathsf N_0},
\]
where \(\mathsf E_0\) is the explicitly constructed rational interval-arithmetic engine supplied by \cref{obj:lem:concrete-arithmetic-engine}, and \(\mathsf N_0\) is the dyadic-floor policy for selecting numerical representations in \cref{obj:def:dyadic-name-convention}.

Each interval is the literal output of the same fixed-budget expression tree, evaluated by its supplied fixed engine on its supplied valid numerical representations. Each choice of engine and representations determines its own rational output, subject to the endpoint enclosure bounds in \cref{obj:lem:fixed-budget-realization}. The construction uses the observable numerator and denominator coordinates \(C(P),S(P)\) from \cref{obj:lem:observable-odds}, separate projection resolutions from \cref{obj:def:resolutions}, and joint inversion from \cref{obj:def:upper-handle}.

The represented inputs are \((n,\alpha,\beta,\mathcal O)\), with numerical representations selected by \(\mathsf N\) and an independently fixed engine \(\mathsf E\) as in \cref{obj:def:arithmetic-engine}. The radius \(r\) indexes evaluation of the expected-length profile.

The converse uses the continuous calibrated priors in \cref{obj:def:continuous-calibrated-priors}, with admissible native propensity and prognosis perturbations, one common scalar effect per draw, analytic singleton matching, and factorization across disjoint latent-sign components conditional on the original covariates. It applies independently of the engine and naming policy.

The model and honesty conditions are specified in \cref{obj:synth_1,obj:synth_2,obj:synth_3,obj:ass:uniform-design,obj:ass:homogeneous-logit,obj:ass:effect-envelope,obj:ass:propensity-envelope,obj:ass:prognosis-envelope,obj:ass:propensity-holder,obj:ass:prognosis-holder,obj:ass:evaluation-radius,obj:ass:global-coverage,obj:def:logistic-class,obj:def:model,obj:def:radius-model,obj:def:honest-procedures,obj:def:length-objective}. The associated constructions are given in \cref{obj:def:causal-extension,obj:def:projection-statistic,obj:def:coordinate-estimates,obj:def:calibration-maps}. Their bounds, calibrations, and sharp comparisons are stated in \cref{obj:lem:projection-control,obj:thm:finite-honest-upper,obj:lem:parametric-null-floor,obj:thm:sharp-honest-frontier,obj:lem:scalar-implicit-function,obj:lem:exact-calibrations,obj:lem:calibrated-support,obj:lem:original-record-mixtures,obj:thm:full-honest-frontier-solution,obj:thm:full-honest-frontier-answer}.
\end{definitionv}

The evaluation radius indexes the expected-length criterion, while the procedure uses the observed sample and public exponents. Thus the same interval sequence can attain the comparison over all radius slices. The following result combines concrete attainment, the lower bound over the entire honest class, and the supporting guarantee for every admissible engine and naming policy.

\begin{theoremv}{thm:full-honest-frontier-answer}[The full honest frontier]
There exist constants \(c,C\in\mathbb R\), with \(c>0\) and \(C>0\), satisfying the following guarantees simultaneously. The public exponent domain and the sharp expected-length profile are
\[
\mathcal D=\{(\alpha,\beta)\in\mathbb R^2:0<\beta<1/4,\ \beta<\alpha\le1\},
\]
\[
\rho_n(\alpha,\beta;r)=d_n(\alpha,\beta;r)
=n^{-\min\{1/2,\,2(\alpha+\beta)/(2\alpha+2\beta+1)\}}
+r\,n^{-4\beta/(4\beta+1)}.
\]

\begin{itemize}
\item \textbf{(Concrete attainment.)}
For every \((\alpha,\beta)\in\mathcal D\) and every integer \(n\ge1\), the concrete interval
\[
I_n^\star(\alpha,\beta)
=I^{\mathrm{up}}_{n,\mathsf E_0,\mathsf N_0}(\alpha,\beta),
\]
constructed with the explicit rational arithmetic engine \(\mathsf E_0\) and the dyadic-floor naming policy \(\mathsf N_0\), belongs to the globally honest class \(\mathcal A_n(\alpha,\beta)\) of \cref{obj:def:honest-procedures}. For every sample and randomization seed, its output is a closed interval \([l,u]\) with
\[
l,u\in\mathbb Q,\qquad -1/2\le l\le u\le1/2.
\]
For every \(r\in[0,1/2]\), the minimax expected length \(J_n(\alpha,\beta;r)\) of \cref{obj:def:length-objective} satisfies
\[
c\,\rho_n(\alpha,\beta;r)
\le J_n(\alpha,\beta;r)
\le
\sup_{P\in\mathcal M_r(\alpha,\beta)}
E_{P^{\otimes n}\otimes\lambda}
\bigl|I_n^\star(\alpha,\beta)(\mathcal O,U)\bigr|
\le C\,\rho_n(\alpha,\beta;r),
\]
where \(\mathcal M_r(\alpha,\beta)\) is the prescribed radius slice and \(\lambda\) is the uniform law of the independent randomization seed.

\item \textbf{(Universal construction.)}
For every admissible rational interval-arithmetic engine \(\mathsf E\), every admissible Borel naming policy \(\mathsf N\), both satisfying the quantitative contracts of \cref{obj:def:arithmetic-engine,obj:def:dyadic-name-convention}, every \((\alpha,\beta)\in\mathcal D\), and every integer \(n\ge1\), the fixed-budget output
\(I^{\mathrm{up}}_{n,\mathsf E,\mathsf N}(\alpha,\beta)\)
is Borel, globally \(90\%\)-honest, and radius independent. On every valid named input, the construction terminates within its prescribed fixed budgets and returns a nonempty closed interval with ordered rational endpoints in \([-1/2,1/2]\), using the observed data and public inputs. For every \(r\in[0,1/2]\),
\[
c\,\rho_n(\alpha,\beta;r)
\le J_n(\alpha,\beta;r)
\le
\sup_{P\in\mathcal M_r(\alpha,\beta)}
E_{P^{\otimes n}\otimes\lambda}
\bigl|I^{\mathrm{up}}_{n,\mathsf E,\mathsf N}(\alpha,\beta)(\mathcal O,U)\bigr|
\le C\,\rho_n(\alpha,\beta;r).
\]
The same constants \(c,C\) apply to all these engines, policies, exponents, sample sizes, and radii.

\item \textbf{(Compact uniformity.)}
Taking the comparison constants \(c_\star=c\), \(C_\star=C\), and the sample-size threshold \(n_0=1\), every nonempty compact set \(K\subseteq\mathcal D\) satisfies
\[
0<\inf_{(\alpha,\beta)\in K}c_\star,\qquad
\sup_{(\alpha,\beta)\in K}C_\star<\infty,\qquad
\sup_{(\alpha,\beta)\in K}n_0=1.
\]

\item \textbf{(Elbow identities.)}
Define the numerator and denominator exponents by
\[
\mathsf a(\alpha,\beta)
=\min\{1/2,\,2(\alpha+\beta)/(2\alpha+2\beta+1)\},
\qquad
\mathsf b(\beta)=\frac{4\beta}{4\beta+1}.
\]
For every \((\alpha,\beta)\in\mathcal D\),
\[
\mathsf a(\alpha,\beta)-\mathsf b(\beta)>0,
\]
and
\[
\mathsf a(\alpha,\beta)-\mathsf b(\beta)
=
\begin{cases}
\displaystyle
\frac{2(\alpha-\beta)}
{(2\alpha+2\beta+1)(4\beta+1)},
&\alpha+\beta\le1/2,\\[6pt]
\displaystyle
\frac{1-4\beta}{2(4\beta+1)},
&\alpha+\beta\ge1/2.
\end{cases}
\]
At \(\alpha+\beta=1/2\), the two displayed expressions are equal.

\item \textbf{(Universal converse.)}
For every \((\alpha,\beta)\in\mathcal D\), every \(r\in[0,1/2]\), every integer \(n\ge1\), and every procedure \(I\in\mathcal A_n(\alpha,\beta)\), including procedures using independent randomization,
\[
c\,\rho_n(\alpha,\beta;r)
\le
\sup_{P\in\mathcal M_r(\alpha,\beta)}
E_{P^{\otimes n}\otimes\lambda}
\bigl|I(\mathcal O,U)\bigr|.
\]
\end{itemize}
\end{theoremv}

The coordinate identity behind the comparison is
\[
C(P)=\{\exp(\theta(P))-1\}S(P),
\qquad
\theta(P)=\log\{1+C(P)/S(P)\},
\qquad S(P)\ge1/512,
\]
as proved in \cref{obj:lem:observable-odds}. Numerator uncertainty therefore contributes directly to interval length, while denominator uncertainty is scaled by the effect radius. Separate projection resolutions accommodate these distinct sources of uncertainty. The matching lower comparison establishes that their sum describes the best possible worst-case expected length under ambient honesty, including for randomized procedures.

The converse has one experiment for each term. In the mixed family, propensity
and prognosis amplitudes \(\eta,\zeta\) produce an effect separation at their
product scale; singleton matching removes first-order marginal evidence, and
the original-record Hellinger bound determines the admissible projection
resolution. This yields the numerator contribution. In the fair-treatment
family, a prognosis amplitude \(\delta\) changes a comparator effect by order
\(t\delta^2\) at center \(t\), while its Hellinger distance is controlled at
fourth order in \(\delta\); taking \(t\) at the evaluation-radius scale yields
the radius-weighted denominator contribution. The continuous sign field shares
endpoint signs across neighboring cells. After conditioning on the observed
covariates, the sign-incidence graph is therefore factored into connected
components and controlled by a spanning-tree bound. The exact matching,
same-model support, and resulting mixture bounds are stated in
\cref{obj:lem:exact-calibrations,obj:lem:calibrated-support,obj:lem:original-record-mixtures};
\cref{obj:lem:parametric-null-floor,obj:thm:sharp-honest-frontier} complete the
two-term comparison.

At the null radius, \cref{obj:thm:full-honest-frontier-answer} gives
\[
c\,n^{-\mathsf a(\alpha,\beta)}
\le J_n(\alpha,\beta;0)
\le C\,n^{-\mathsf a(\alpha,\beta)}.
\]
Coverage remains imposed over the ambient model through \cref{obj:def:honest-procedures}, even though expected length is evaluated at zero effect. The null comparison therefore measures the precision attainable at zero by a procedure that also covers every effect in the prescribed ambient envelope. Its rate is governed by the numerator contribution.

For any sequence \(r_n\in[0,1/2]\), including an arbitrary shrinking-radius sequence, the simultaneous comparison gives
\[
c\left\{
n^{-\mathsf a(\alpha,\beta)}
+r_n n^{-\mathsf b(\beta)}
\right\}
\le J_n(\alpha,\beta;r_n)
\le
C\left\{
n^{-\mathsf a(\alpha,\beta)}
+r_n n^{-\mathsf b(\beta)}
\right\}.
\]
The same concrete interval sequence attains these bounds across the sequence. Equating the two contributions yields the radius balance
\[
r_n=n^{-\{\mathsf a(\alpha,\beta)-\mathsf b(\beta)\}}.
\]
For fixed public exponents, radii smaller than this scale make the numerator contribution dominant; radii larger than this scale make the denominator contribution dominant. At the balance scale, both contributions have the same order.

The numerator transition occurs at \(\alpha+\beta=1/2\). Below this boundary, its exponent is \(2(\alpha+\beta)/(2\alpha+2\beta+1)\), and the radius-balance exponent depends on the difference \(\alpha-\beta\). At and above the boundary, the numerator exponent equals \(1/2\), and the balance exponent depends on \(\beta\) alone. The two formulas agree at the shared boundary, so the radius balance is continuous across the numerator transition.

The propensity endpoint \(\alpha=1\) belongs to the public exponent domain in \cref{obj:thm:full-honest-frontier-answer}. At this endpoint, every admissible prognosis exponent places the numerator in the parametric regime, with balance exponent
\[
\mathsf a(1,\beta)-\mathsf b(\beta)
=
\frac{1-4\beta}{2(4\beta+1)}.
\]
Along the prognosis boundary approached from \(\beta<1/4\), this expression tends to zero, while the denominator exponent tends to \(1/2\). Within the stated domain, the positive exponent gap ensures that for every fixed \(r\in(0,1/2]\), the radius-scaled denominator contribution eventually dominates. Consequently, the sharp expected length at a fixed positive radius has order \(r n^{-\mathsf b(\beta)}\).

Concrete attainment is a guarantee for a statistical interval procedure with rational endpoints on every sample. The engine and naming-policy clause attaches the same coverage and expected-length comparison to every admissible numerical realization, with common constants. The following construction specifies the interval itself in \cref{obj:def:upper-handle,obj:thm:finite-honest-upper}, and \cref{obj:lem:concrete-arithmetic-engine,obj:lem:fixed-budget-realization} supplies its arithmetic realization.
